%% ============================================================
%% Chapter 2 Solutions
%% ============================================================

\section{Chapter 2: The Sensitivity Index}

%% ----------------------------------------------------------
\begin{solution}{2.1}{L2}

\solanswer

The sign of a sensitivity index tells you the \emph{direction} of the
relationship: a positive SI means the \QOI increases when the \POI
increases; a negative SI means the \QOI decreases when the \POI increases.
The sign is the sign of the underlying derivative $\partial q/\partial p$,
scaled by $p/q$ (both positive under the interpretable-parameterization
convention).

\solpart{a} \emph{Lung capacity vs.\ cigarettes per day.}
Lung capacity decreases as cigarette consumption increases (medical
consensus). Therefore $\partial q/\partial p < 0$, so the SI is
\textbf{negative}.

\solpart{b} \emph{Crop yield vs.\ annual rainfall (near the optimum).}
Crop yields typically increase with rainfall up to an optimum and
then decrease (waterlogging). The sign depends on which side of the
optimum we are evaluating. Near the dry end of the feasible range,
yield increases with rainfall: \textbf{positive} SI. Near the wet
end (waterlogging), yield decreases: \textbf{negative} SI. The question
specifies ``near a region of the parameter space''; the sign must be
specified relative to that region, and the answer should note the
non-monotonicity.

\solpart{c} \emph{$R_0$ vs.\ $\gamma$ (recovery rate) for SIR.}
$R_0 = k\beta/\gamma$, so $\partial R_0/\partial\gamma < 0$: faster
recovery reduces $R_0$. The SI is \textbf{negative}: $\SI{\gamma}{R_0} = -1$.

\solnote{Part (b) is deliberately ambiguous to prompt students to
identify the dependence on the evaluation point. Accept either sign
if the student correctly identifies the nonmonotonicity and states
which side of the optimum they are considering. This previews Chapter 8
(when SA breaks down near parameter thresholds).}

\end{solution}

%% ----------------------------------------------------------
\begin{solution}{2.2}{L2}

\solanswer

Let $q = \sqrt{p}$, so $q = p^{1/2}$.

\textbf{Form 1 (limit definition):}
\[
  \SI{p}{q} = \lim_{\delta p \to 0} \frac{\delta q/q}{\delta p/p}
  = \lim_{\delta p \to 0}
    \frac{\bigl(\sqrt{p+\delta p}-\sqrt{p}\bigr)/\sqrt{p}}
         {\delta p/p}
  = \frac{p}{\sqrt{p}} \lim_{\delta p\to 0}
    \frac{\sqrt{p+\delta p}-\sqrt{p}}{\delta p}
  = \sqrt{p}\cdot\frac{1}{2\sqrt{p}} = \frac{1}{2}.
\]

\textbf{Form 2 (derivative form):}
\[
  \SI{p}{q} = \frac{p}{q}\,\frac{\partial q}{\partial p}
  = \frac{p}{\sqrt{p}}\cdot\frac{1}{2\sqrt{p}}
  = \frac{\sqrt{p}}{2\sqrt{p}} = \frac{1}{2}.
\]

\textbf{Form 3 (logarithmic form):}
\[
  \SI{p}{q} = \frac{d\ln q}{d\ln p}
  = \frac{d}{d\ln p}\bigl(\tfrac{1}{2}\ln p\bigr)
  = \frac{1}{2}.
\]

All three give $\SI{p}{q} = \tfrac{1}{2}$.

\textbf{Interpretation:} A $1\%$ increase in $p$ produces approximately
a $0.5\%$ increase in $q = \sqrt{p}$. More precisely: if $p$ doubles
(increases by $100\%$), $q$ increases by approximately $50\%$ (the exact
value is $\sqrt{2}-1 \approx 41.4\%$; the SI gives a linear approximation).

\solnote{Emphasize the logarithmic form (Form 3) as the most computationally
efficient: $d\ln q/d\ln p$ can often be read off immediately from the
algebraic form of $q(p)$. For a power law $q = cp^n$, all three forms
give $\SI{p}{q} = n$, which can be seen immediately from the log form.}

\end{solution}

%% ----------------------------------------------------------
\begin{solution}{2.3}{L3}

\solanswer

Let $\mu = \lambda b$. The arc-length functional is
\[
  J = \int_0^b \sqrt{1+\lambda^2 t^2}\,dt.
\]

\textbf{Evaluating $J$.} Substitute $u = \lambda t$, $du = \lambda\,dt$:
\[
  J = \frac{1}{\lambda}\int_0^{\lambda b}\sqrt{1+u^2}\,du
    = \frac{1}{\lambda}\Bigl[\tfrac{u}{2}\sqrt{1+u^2}
      + \tfrac{1}{2}\arcsinh u\Bigr]_0^{\mu}
    = \frac{\mu\sqrt{1+\mu^2} + \arcsinh\mu}{2\lambda}.
\]

\textbf{Computing $\SI{b}{J}$.} Using the derivative form with $\mu = \lambda b$:
\[
  \frac{\partial J}{\partial b} = \sqrt{1+\lambda^2 b^2} = \sqrt{1+\mu^2}.
\]
So
\[
  \SI{b}{J} = \frac{b}{J}\cdot\sqrt{1+\mu^2}
  = \frac{\lambda b\cdot\sqrt{1+\mu^2}}{\tfrac{1}{2}(\mu\sqrt{1+\mu^2}+\arcsinh\mu)}
  = \frac{2\mu\sqrt{1+\mu^2}}{\mu\sqrt{1+\mu^2}+\arcsinh\mu}.
\]
Factor $\sqrt{1+\mu^2}$ from numerator and denominator:
\[
  \SI{b}{J} = \frac{2\mu}{\mu + \arcsinh\mu/\sqrt{1+\mu^2}}.
\]
This matches the stated formula. $\checkmark$

\textbf{Computing $\SI{\lambda}{J}$.} Differentiating $J = (\mu\sqrt{1+\mu^2}+\arcsinh\mu)/(2\lambda)$
with $\mu = \lambda b$ (so $\partial\mu/\partial\lambda = b = \mu/\lambda$):

\[
  \frac{\partial J}{\partial\lambda}
  = \frac{b\bigl(\sqrt{1+\mu^2} + \mu^2/\sqrt{1+\mu^2} + 1/\sqrt{1+\mu^2}\bigr)}{2\lambda}
  - \frac{J}{\lambda}.
\]
Simplifying $\sqrt{1+\mu^2} + \mu^2/\sqrt{1+\mu^2} + 1/\sqrt{1+\mu^2}
= (1+\mu^2+\mu^2+1)/\sqrt{1+\mu^2} = (2+2\mu^2)/\sqrt{1+\mu^2} = 2\sqrt{1+\mu^2}$:
\[
  \frac{\partial J}{\partial\lambda}
  = \frac{\mu\sqrt{1+\mu^2}}{\lambda^2} - \frac{J}{\lambda}.
\]
Therefore
\[
  \SI{\lambda}{J}
  = \frac{\lambda}{J}\cdot\frac{\partial J}{\partial\lambda}
  = \frac{\lambda}{J}\!\left(\frac{\mu\sqrt{1+\mu^2}}{\lambda^2} - \frac{J}{\lambda}\right)
  = \frac{\mu\sqrt{1+\mu^2}}{\lambda J} - 1.
\]
Substituting $J = (\mu\sqrt{1+\mu^2}+\arcsinh\mu)/(2\lambda)$:
\[
  \SI{\lambda}{J}
  = \frac{2\mu\sqrt{1+\mu^2}}{\mu\sqrt{1+\mu^2}+\arcsinh\mu} - 1
  = \frac{2\mu\sqrt{1+\mu^2} - \mu\sqrt{1+\mu^2} - \arcsinh\mu}
         {\mu\sqrt{1+\mu^2}+\arcsinh\mu}
  = \frac{\mu\sqrt{1+\mu^2} - \arcsinh\mu}
         {\mu\sqrt{1+\mu^2}+\arcsinh\mu}.
\]
This matches the stated formula. $\checkmark$

\solnote{This is primarily an integration and chain-rule exercise.
The key pedagogical point: both SIs depend only on $\mu = \lambda b$
(the dimensionless parameter group), illustrating why dimensional
analysis and natural parameterization matter. Full credit requires
correct derivation of both formulas; partial credit for correct
setup with algebra error.}

\end{solution}

%% ----------------------------------------------------------
\begin{solution}{2.4}{L3}

\solanswer

Let $J_+ = (-b + \sqrt{b^2-4ac})/(2a)$. Denote $D = b^2-4ac > 0$.
Note that $J_+$ satisfies $aJ_+^2 + bJ_+ + c = 0$, so $aJ_+^2 = -bJ_+ - c$.

The denominator in all three formulas is
$2aJ_+ + b = \sqrt{b^2-4ac} = \sqrt{D}$
(from the quadratic formula: $2aJ_+ = -b + \sqrt{D}$, so $2aJ_+ + b = \sqrt{D}$).

\textbf{Computing $\SI{a}{J_+}$:}
\[
  \frac{\partial J_+}{\partial a}
  = \frac{-b+\sqrt{D}}{-2a^2} + \frac{-4c/2}{2a\sqrt{D}}
  \cdot(-1)
  \quad\text{[quotient rule + chain rule]}
\]
Simpler via implicit differentiation: differentiate $aJ_+^2 + bJ_+ + c = 0$
with respect to $a$:
\[
  J_+^2 + (2aJ_+ + b)\frac{\partial J_+}{\partial a} = 0
  \implies
  \frac{\partial J_+}{\partial a} = -\frac{J_+^2}{2aJ_++b}.
\]
Therefore
\[
  \SI{a}{J_+} = \frac{a}{J_+}\cdot\frac{\partial J_+}{\partial a}
  = \frac{a}{J_+}\cdot\frac{-J_+^2}{2aJ_++b}
  = \frac{-aJ_+}{2aJ_++b}. \checkmark
\]

\textbf{Computing $\SI{b}{J_+}$:}
Differentiate $aJ_+^2 + bJ_+ + c = 0$ with respect to $b$:
\[
  J_+ + (2aJ_++b)\frac{\partial J_+}{\partial b} = 0
  \implies
  \frac{\partial J_+}{\partial b} = -\frac{J_+}{2aJ_++b}.
\]
\[
  \SI{b}{J_+} = \frac{b}{J_+}\cdot\frac{-J_+}{2aJ_++b}
  = \frac{-b}{2aJ_++b}. \checkmark
\]

\textbf{Computing $\SI{c}{J_+}$:}
Differentiate with respect to $c$:
\[
  (2aJ_++b)\frac{\partial J_+}{\partial c} + 1 = 0
  \implies
  \frac{\partial J_+}{\partial c} = -\frac{1}{2aJ_++b}.
\]
\[
  \SI{c}{J_+} = \frac{c}{J_+}\cdot\frac{-1}{2aJ_++b}
  = \frac{-c}{J_+(2aJ_++b)}. \checkmark
\]

\textbf{Numerical verification at $a=1$, $b=-3$, $c=2$:}
$J_+ = (3+\sqrt{9-8})/2 = 2$, $2aJ_++b = 4-3 = 1$.
\[
  \SI{a}{J_+} = -\frac{1\cdot 2}{1} = -2,\quad
  \SI{b}{J_+} = -\frac{-3}{1} = 3,\quad
  \SI{c}{J_+} = -\frac{2}{2\cdot 1} = -1.
\]
Verification by finite differences (centered, $h = 10^{-6}$):
$\partial J_+/\partial a \approx (J_+(a+h)-J_+(a-h))/(2h)$,
multiplied by $a/J_+ = 1/2$, gives $\SI{a}{J_+} \approx -2.000000$ $\checkmark$.

\solnote{Implicit differentiation is the most elegant approach and avoids
messy quotient-rule algebra. Students who use explicit differentiation
of the quadratic formula are not wrong but often make sign errors.
The numerical check at $a=1$, $b=-3$, $c=2$ is important: students
should verify that their symbolic answers match finite differences.}

\end{solution}

%% ----------------------------------------------------------
\begin{solution}{2.5}{L3}

\solanswer

The SEIR reproduction number is
\[
  R_0 = \frac{k\beta\nu}{(\nu+\mu)(\gamma+\mu)}.
\]
Reparameterize: $\tau_E = 1/\nu$, $\tau_I = 1/\gamma$, $L = 1/\mu$.
Substituting $\nu = 1/\tau_E$, $\gamma = 1/\tau_I$, $\mu = 1/L$:

\[
  R_0 = \frac{k\beta/\tau_E}{(1/\tau_E + 1/L)(1/\tau_I + 1/L)}
  = \frac{k\beta L^2\tau_I}{(\tau_E + L)(\tau_I\tau_E + \tau_E L + \tau_I L + L^2/\tau_E)}
\]
For $L \gg \tau_E, \tau_I$ (lifespan much longer than disease timescales,
which is $10{,}000 \gg 7$ in the SIR examples):
\[
  R_0 \approx k\beta\tau_I\tau_E\cdot\frac{\nu}{\nu+\mu}\cdot\frac{1}{1+\mu/\gamma}
  \approx k\beta\tau_I \cdot \frac{\tau_E}{1}
  \quad\text{(since } \mu\ll\nu, \mu\ll\gamma).
\]

\textbf{Sensitivity indices} (using log-derivative form):
Writing $R_0 = k\beta\nu/[(\nu+\mu)(\gamma+\mu)]$:
\[
  \ln R_0 = \ln k + \ln\beta + \ln\nu - \ln(\nu+\mu) - \ln(\gamma+\mu).
\]
\[
  \SI{k}{R_0} = 1, \qquad \SI{\beta}{R_0} = 1.
\]
\[
  \SI{\nu}{R_0} = \frac{d\ln R_0}{d\ln\nu}
  = 1 - \frac{\nu}{\nu+\mu} = \frac{\mu}{\nu+\mu}.
\]
In duration parameterization ($\tau_E = 1/\nu$, and $\SI{\tau_E}{R_0} = -\SI{\nu}{R_0}$):
\[
  \SI{\tau_E}{R_0} = -\frac{\mu}{\nu+\mu} = -\frac{1/L}{1/\tau_E+1/L}
  = \frac{-\tau_E}{L+\tau_E}.
\]
Similarly, $\SI{\gamma}{R_0} = -\mu/(\gamma+\mu)$ and
\[
  \SI{\tau_I}{R_0} = \frac{\mu}{\gamma+\mu}/({\text{sign flip}})
  = \frac{\mu}{\gamma+\mu}\cdot\frac{\tau_I}{\tau_I}
  \implies
  \SI{\tau_I}{R_0} = \frac{1}{1+\mu\tau_I} = \frac{L}{L+\tau_I}.
\]
For nominal values $k=5$, $\beta=0.06$, $\tau_E=5.5$ days,
$\tau_I=7$ days, $L=10{,}000$ days:
\begin{center}
\begin{tabular}{lcccc}
\toprule
\POI & $k$ & $\beta$ & $\tau_E$ & $\tau_I$ \\
\midrule
SI   & $1.000$ & $1.000$ & $-0.00055$ & $0.0007$ \\
\bottomrule
\end{tabular}
\end{center}
For $L \gg \tau_E, \tau_I$: $\SI{\tau_E}{R_0} \approx -\tau_E/L \approx 0$
and $\SI{\tau_I}{R_0} \approx \tau_I/L \approx 0$.
In the demographic limit, $k$ and $\beta$ dominate, and both $\tau_E$
and $\tau_I$ have negligible SI because the mortality rate is so small.
(If $L$ were short, e.g., $L = 30$ days, the disease timescales
would matter significantly.)

\solnote{Key lesson: in models with demography, the lifespan $L$ modulates
how much the exposed period matters. Students often expect $\tau_E$ to
have a larger SI; explaining why it doesn't requires understanding
$\mu/(\nu+\mu) = \tau_E/(\tau_E+L)$.}

\end{solution}

%% ----------------------------------------------------------
\begin{solution}{2.6}{L4}

\solanswer

\textbf{Paragraph for the policy maker:}

Both papers are reporting the same underlying biology, but in different
parameterizations. The first paper uses $\gamma$, the recovery rate
(units: $\text{day}^{-1}$), while the second uses $\tau = 1/\gamma$,
the mean infectious period (units: days). For the SIR model,
$R_0 = k\beta/\gamma = k\beta\tau$, so increasing $\gamma$
decreases $R_0$ (negative index), while increasing $\tau$ increases
$R_0$ (positive index). The reciprocal relationship $\tau = 1/\gamma$
means that ``$\gamma$ increases'' and ``$\tau$ decreases'' describe
the same biological event: faster recovery. Both papers are correct;
the sign difference reflects the direction of the parameterization,
not a scientific disagreement.

For policy, the duration parameterization ($\tau$) is more informative:
public health officers think in terms of ``the average patient is
infectious for 7 days,'' not ``the recovery rate is 0.14 per day.''
A positive $\SI{\tau}{R_0} = +1$ communicates directly that reducing
the infectious period by $10\%$ (through treatment or isolation) would
reduce $R_0$ by $10\%$. Reporting should use the duration parameterization
and state explicitly that reducing $\tau$ reduces $R_0$.

\solnote{The one-paragraph format is deliberate: the exercise tests
whether students can synthesize the parameterization concept into a
clear, jargon-free explanation. Deduct credit if the student simply
writes equations without a clear explanation for a non-technical reader.
The policy-relevance component (which parameterization is more actionable)
is essential for full credit.}

\end{solution}

%% ----------------------------------------------------------
\begin{solution}{2.7}{L4}

\solanswer

The sensitivity indices (sorted by $|\cdot|$, descending):
\begin{center}
\begin{tabular}{lcc}
\toprule
\POI & SI & $|\text{SI}|$ \\
\midrule
$\tau_c$ (contamination duration) & $+2.1$ & $2.1$ \\
$\beta_f$ (food-to-person transmission) & $+1.8$ & $1.8$ \\
$\tau_s$ (hand-washing interval) & $-1.3$ & $1.3$ \\
$k_f$ (food contacts per day) & $+0.7$ & $0.7$ \\
$\mu$ (clearance rate) & $-0.2$ & $0.2$ \\
\bottomrule
\end{tabular}
\end{center}

\solpart{a} MATLAB code using \texttt{tornado\_plot} from the repository:
\begin{verbatim}
SI     = [2.1; 1.8; -1.3; 0.7; -0.2];
labels = {'$\tau_c$','$\beta_f$','$\tau_s$','$k_f$','$\mu$'};
tornado_plot(SI, labels, 'J');   %% J = mean illness cases
\end{verbatim}
The resulting plot has $\tau_c$ as the widest bar (top), extending
right to $+2.1$; $\mu$ as the shortest bar (bottom), extending left
to $-0.2$.

\solpart{b} \textbf{Intervention priority:}
Target $\tau_c$ first (SI $= +2.1$): reducing contamination duration
has the largest effect on illness cases. Shortening the contamination
event (e.g., faster detection and decontamination protocols) would
reduce cases roughly proportionally.

\solpart{c} \textbf{Interpreting $\SI{\tau_c}{J} = +2.1$:}
A $1\%$ reduction in mean contamination duration reduces illness
cases by approximately $2.1\%$. The index exceeds 1 because the
model is nonlinear in $\tau_c$ (total cases grow super-linearly with
contamination duration; longer events expose more people per event).

\solpart{d} \textbf{$\SI{\tau_s}{J} = -1.3$ (negative):}
More frequent hand-washing (smaller $\tau_s$ = shorter interval between
washings) reduces transmission, so illness cases decrease as $\tau_s$
decreases. ``Increasing $\tau_s$'' means less frequent hand-washing,
which increases cases. Hence the positive-$\tau_s$ direction corresponds
to increased illness: SI is negative (longer interval between washings $\to$
more transmission). The intervention is to \emph{decrease} $\tau_s$
(increase hand-washing frequency).

\solnote{Part (c) is a good teaching moment: SI $> 1$ does not
indicate an error. For a model that is super-linear in a parameter
(e.g., $J \propto \tau^2$), the SI is $2$, not $1$. This is one
of the cases where the sign and magnitude carry richer information
than a simple linear model would predict.}

\end{solution}

%% ----------------------------------------------------------
\begin{solution}{2.8}{L5}

\solanswer

\textbf{Model answer (design for a six-parameter respiratory model).}

\solpart{a} \textbf{Parameterization.}
Use duration parameterization throughout: mean latent period $\tau_E$,
mean infectious period $\tau_I$, mean immunity duration $\tau_R$
(if applicable), and mean lifespan $L$ in days. Report contact
rates $k$ (contacts/person/day) and transmission probability $\beta$
as dimensionless or rate parameters but interpret sensitivity indices
in terms of ``a 10-day increase in $\tau_I$'' rather than
``a 0.14 decrease in $\gamma$.'' Rationale: the WHO team includes
epidemiologists and clinicians; duration units are actionable
(isolation protocols target $\tau_I$; vaccine protection duration
targets $\tau_R$).

\solpart{b} \textbf{QOI selection.}
Primary: total deaths at 6 months (directly relevant to allocation decisions).
Secondary: $R_0$ (determines whether epidemic occurs at all, a natural
threshold). Optional: peak hospitalization rate (determines whether
healthcare capacity is overwhelmed). Using multiple \QOIs catches
cases where the most important parameter for $R_0$ is different from
the most important parameter for peak hospitalization.

\solpart{c} \textbf{Method.}
Start with local SA (forward sensitivity equations or centered finite
differences, 6 parameters, $\approx 13$ model evaluations). This gives
fast rankings to guide early decisions. If the model is nonlinear
(identifiable only from aggregate data, or exhibiting threshold
behavior near $R_0 = 1$), follow with Sobol indices on the top
3--4 parameters. If each evaluation takes seconds, a full Sobol at
$N = 1000$ is feasible immediately.

\solpart{d} \textbf{Validation.}
Check that SI values for $R_0$ match the analytic values ($\SI{k}{R_0} = \SI{\beta}{R_0} = \SI{\tau_I}{R_0} = 1$) using the analytic formula. Use this as a sanity check before reporting any sensitivity results for deaths.

\solpart{e} \textbf{Limitation.}
Local SA is evaluated at a single fitted parameter set. For a new
pathogen, the true parameters may be far from the fitted values.
The SI rankings might change substantially across the plausible
parameter range. I would report this limitation explicitly and
flag any parameter whose SI changes sign or magnitude substantially
across a 50\% confidence interval around the fitted values.

\solnote{Rubric: 1 point each for (a) justification of parameterization,
(b) motivated \QOI choice, (c) specific SA method with cost estimate,
(d) validation procedure, (e) genuine limitation. Five points total.}

\end{solution}

%% ----------------------------------------------------------
\begin{solution}{2.9}{L6}

\solanswer

\textbf{The error:} The recovery rate $\gamma$, mortality rate $\mu$,
and waning immunity rate $\rho$ are all \emph{rates} (units: $\text{day}^{-1}$);
increasing them corresponds to faster recovery, faster death, and
faster waning, respectively. For standard epidemic models:
\begin{itemize}
  \item $\SI{\gamma}{R_0} = -1$ (faster recovery reduces $R_0$) --- \textbf{negative}.
  \item $\SI{\mu}{R_0}$ is typically negative (background mortality
        reduces effective infectious period).
  \item $\SI{\rho}{R_0}$ for a model with waning immunity is also
        typically negative or zero (faster waning increases
        susceptibility, which could increase $R_0$ --- this one
        is genuinely positive in some formulations).
\end{itemize}
The claim that ``all sensitivity indices were positive'' is inconsistent
with these mathematical relationships for at least $\gamma$ and $\mu$.

\textbf{What specifically is wrong:}
The authors either (a) computed derivatives without the normalization
factor $p/q$ (so they computed $\partial R_0/\partial\gamma < 0$,
giving negative sign, but then reported absolute values); or
(b) reparameterized implicitly and are reporting $\SI{\tau_I}{R_0} = +1$
(which is positive) without acknowledging the parameterization switch;
or (c) made a coding error that reversed the sign. Option (b) is
most charitable and most common.

\textbf{What should have been done:}
The paper should state explicitly which parameterization was used
(rates or durations) and report the actual signed SI values. If
duration parameterization was used, the table should identify
each parameter as a rate or duration to avoid confusion. A statement
that ``all sensitivity indices were positive'' without this context
is ambiguous and potentially incorrect.

\solnote{This exercise targets a common real-world error in published
papers. The key diagnostic is: rates and reproduction numbers always
have negative SI relationships (increasing rates reduces $R_0$
through faster depletion of infectious individuals). Students who
only identify ``wrong sign'' without explaining \emph{why} should
receive partial credit. Full credit requires identifying the
mechanism (normalization, parameterization, or sign error).}

\end{solution}
